Sentiment analysis determines the polarity of a piece of text, and product reviews are one of its most common applications.
This report builds and evaluates classifiers that label a product review as negative (0) or positive (1).
The training set contains 7{,}401 labeled reviews, of which 85.4\% are positive, and the test set contains 1{,}851 unlabeled reviews.
The task combines a modest amount of training data, a strong class imbalance and noisy, informal text, so the choice of feature format matters for both accuracy and cost.

We make four contributions.
(1) We define one grouped, stratified validation split in which duplicate review texts never cross the split, and use it for every model, so that all scores are comparable.
(2) We compare three feature formats of increasing cost under identical conditions: sparse TF-IDF n-grams with linear classifiers, static GloVe embeddings with logistic regression and a BiLSTM, and subword tokens with a fine-tuned RoBERTa-base model.
(3) We submit the average of three RoBERTa runs, which reaches 0.962 accuracy and 0.927 Macro-F1 on the validation split, and we analyze its predictions on a manually labeled sample of test reviews.
(4) We discuss how the approach transfers to unlabeled hotel reviews and, with a controlled label-flipping experiment, how noisy ratings degrade it and which of three remedies helps most.

Section~\ref{sec:related} reviews the literature. Sections~\ref{sec:data}--\ref{sec:hparams} describe the data, the models and the parameter settings; Sections~\ref{sec:results}--\ref{sec:formats} report the test-set results, the analysis of predictions and the comparison of feature formats; Sections~\ref{sec:hotel} and~\ref{sec:noise} address the two transfer scenarios, and Section~\ref{sec:conclusion} concludes.
Table~\ref{tab:intro-map} shows where each question of the assignment is answered, and Appendix~\ref{app:code} explains how to reproduce the results.

\begin{table}[tbp]
\centering
\caption{Where each question of the assignment is answered.}
\label{tab:intro-map}
\small
\begin{tabular}{l l l}
\toprule
Question & Topic & Section \\
\midrule
(a) & Literature survey & \ref{sec:related} (\ref{sec:problem}--\ref{sec:baseline}) \\
(b) & Feature format and data preprocessing & \ref{sec:data} \\
(c) & Model, architecture, loss and training & \ref{sec:models} \\
(d) & How parameters were determined & \ref{sec:hparams} \\
(e) & Test-set results and \texttt{submission.csv} & \ref{sec:results} \\
(f) & Correctly and incorrectly classified samples & \ref{sec:analysis} \\
(g) & Feature format, resources and accuracy & \ref{sec:formats} \\
(h) & Hotel reviews without rating scores & \ref{sec:hotel} \\
(i) & Hotel reviews with noisy rating scores & \ref{sec:noise} \\
\bottomrule
\end{tabular}
\end{table}
